% =============================================================================
% Worked solutions for Simulados/mock3-midterm-topics-1-2b.pdf (Topics 1, 2A, 2B).
% Statements: mock3-midterm-topics-1-2b-statements.tex (shared with the exam paper).
% Compile: python -I Simulados/mock3-midterm-topics-1-2b-solutions_check.py (must print PASS),
%          then pdflatex (twice); check: pdffonts (Type 1 only, uni = yes).
% =============================================================================
\documentclass[11pt,a4paper]{article}
\usepackage[utf8]{inputenc}
\usepackage[T1]{fontenc}
\usepackage{lmodern}   % Type 1 vector fonts (replaces PK bitmaps)
\usepackage{cmap}      % ToUnicode CMap -> copyable text and accents
\usepackage[english]{babel}
\usepackage{amsmath,amssymb,amsfonts}
\usepackage[a4paper,margin=2.2cm]{geometry}
\usepackage{booktabs}
\usepackage{enumitem}
\usepackage{fancyhdr}
\usepackage{xcolor}
\usepackage{microtype}

\pagestyle{fancy}\fancyhf{}
\lhead{\footnotesize Finan\c{c}as II -- Corporativas (MPE Insper)}
\rhead{\footnotesize Mock midterm 3, Topics 1--2B --- Worked solutions}
\cfoot{\footnotesize\thepage}
\renewcommand{\headrulewidth}{0.3pt}
\setlength{\headheight}{14pt}

\makeatletter
\renewcommand{\@seccntformat}[1]{\csname the#1\endcsname.\quad}
\makeatother
\numberwithin{equation}{section}
\setlist{topsep=2pt,itemsep=1pt}

\newcommand{\srcline}[1]{\par\noindent{\small\color{black!65}\textbf{Source:} #1}\par\smallskip}
\newcommand{\statement}[1]{\par\smallskip\noindent\colorbox{black!5}{\parbox{\dimexpr\linewidth-2\fboxsep}{\small #1}}\par\medskip}
\newcommand{\step}[2]{\par\medskip\noindent\textbf{Step #1 --- #2.}~}
\newenvironment{notation}{\par\noindent\textbf{Notation.}\begin{itemize}[leftmargin=1.4em]}{\end{itemize}}
\newcommand{\intuition}[1]{\par\medskip\noindent\textbf{Intuition.}~#1}
\newcommand{\refline}[1]{\par\smallskip\noindent{\small\color{black!70}\textbf{Ref:} #1}}
\newcommand{\rubric}[2]{\par\smallskip\noindent{\small\textbf{Marking guide (#1 marks).}~#2}}
\newcommand{\answer}[1]{\par\medskip\noindent\fbox{\parbox{\dimexpr\linewidth-2\fboxsep-2\fboxrule}{\textbf{Answer.}~#1}}\par}
\newcommand{\tf}[3]{\item[#1] \textbf{#2.} #3}

% =============================================================================
% Statements of mock midterm 3 (Topics 1, 2A, 2B), written once and \input by
%   mock3-midterm-topics-1-2b.tex            (exam paper)
%   mock3-midterm-topics-1-2b-solutions.tex  (worked solutions)
% Every statement is a real exercise copied verbatim (rules/exercise-lists.md):
% source line first, then the full text. Allowed edits only: math in LaTeX, tables
% rebuilt, line breaks and hyphenation undone, currency written as \$.
% Needs: \srcline{...} defined by the including document.
% =============================================================================

% ---------------------------------------------------------------- Question 1.1
\newcommand{\QoneA}{%
\srcline{Higgins (Koski--Higgins--Mitton), \emph{Analysis for Financial Management} 13e, Ch.~2,
Exercise 3, pp.~71--72 (PDF 92--93) $\cdot$ Answer in book: yes, p.~401 (PDF 422)}

\noindent Answer true or false to each of the following. Briefly explain your reasoning for each answer.
\begin{enumerate}[label=\alph*.,leftmargin=2em]
  \item A company's assets-to-equity ratio always equals one plus its liabilities-to-equity ratio.
  \item A company's return on equity will always equal or exceed its return on assets.
  \item A company's collection period should always be less than its payables period.
  \item A company's current ratio must always equal or exceed its acid-test ratio.
  \item All else equal, a firm would prefer to have a higher asset turnover ratio.
  \item Two firms can have the same earnings yield but different price-to-earnings ratios.
  \item Ignoring taxes and transactions costs, unrealized paper gains are less valuable than realized
        cash earnings.
\end{enumerate}}

% ---------------------------------------------------------------- Question 1.2
\newcommand{\QoneB}{%
\srcline{Higgins (Koski--Higgins--Mitton), \emph{Analysis for Financial Management} 13e, Ch.~4,
Exercise 3, p.~134 (PDF 155) $\cdot$ Answer in book: yes, pp.~408--409 (PDF 429--430)}

\noindent Are the following statements true or false? Please explain why.
\begin{enumerate}[label=\alph*.,leftmargin=2em]
  \item The only way a company can grow at a rate above its current sustainable growth rate is by
        issuing new stock.
  \item The stock market is a ready source of new capital when a company is incurring heavy losses.
  \item Share repurchases usually increase earnings per share.
  \item Companies often buy back their stock because managers believe the shares are undervalued.
  \item Only rapidly growing firms have growth management problems.
  \item Increasing growth increases stock price.
\end{enumerate}}

% ---------------------------------------------------------------- Question 1.3
\newcommand{\QoneC}{%
\srcline{Berk, DeMarzo \& Harford, \emph{Fundamentals of Corporate Finance} 2e, Ch.~16, Critical
Thinking question 2, p.~490 (PDF 521) $\cdot$ Answer in book: no}

\noindent Explain what is wrong with the following argument: ``If a firm issues debt that is risk-free
because there is no possibility of default, the risk of the firm's equity does not change. Therefore,
risk-free debt allows the firm to get the benefit of a low cost of capital of debt without raising its
cost of capital of equity.''}

% ---------------------------------------------------------------- Question 2 (section instruction)
\newcommand{\QtwoIntro}{%
\noindent\emph{Section instruction} (Berk, Ch.~16, ``Capital Structure in Perfect Capital Markets'',
p.~490, PDF 521, applying to Parts 2.1, 2.2 and 2.4): For problems in this section assume no taxes or
distress costs.}

% ---------------------------------------------------------------- Question 2.1
\newcommand{\QtwoA}{%
\srcline{Berk, DeMarzo \& Harford, \emph{Fundamentals of Corporate Finance} 2e, Ch.~16, Exercise 16.2,
p.~491 (PDF 522) $\cdot$ Answer in book: no}

\noindent You are an entrepreneur starting a biotechnology firm. If your research is successful, the
technology can be sold for \$30 million. If your research is unsuccessful, it will be worth nothing. To
fund your research, you need to raise \$2 million. Investors are willing to provide you with \$2 million
in initial capital in exchange for 50\% of the unlevered equity in the firm.
\begin{enumerate}[label=\alph*.,leftmargin=2em]
  \item What is the total market value of the firm without leverage?
  \item Suppose you borrow \$1 million. According to MM, what fraction of the firm's equity will you
        need to sell to raise the additional \$1 million you need?
  \item What is the value of your share of the firm's equity in cases (a) and (b)?
\end{enumerate}}

% ---------------------------------------------------------------- Question 2.2
\newcommand{\QtwoB}{%
\srcline{Berk, DeMarzo \& Harford, \emph{Fundamentals of Corporate Finance} 2e, Ch.~16, Exercise *16.4,
p.~491 (PDF 522) $\cdot$ Answer in book: no}

\noindent Suppose there are no taxes. Firm ABC has no debt, and firm XYZ has debt of \$5000 on which it
pays interest of 10\% each year. Both companies have identical projects that generate free cash flows of
\$800 or \$1000 each year. After paying any interest on debt, both companies use all remaining free cash
flows to pay dividends each year.
\begin{center}
\begin{tabular}{lcccc}
  \toprule
  & \multicolumn{2}{c}{ABC} & \multicolumn{2}{c}{XYZ} \\
  \cmidrule(lr){2-3}\cmidrule(lr){4-5}
  FCF & Debt Payments & Equity Dividends & Debt Payments & Equity Dividends \\ \midrule
  \$\ 800 & & & & \\
  \$1000 & & & & \\ \bottomrule
\end{tabular}
\end{center}
\begin{enumerate}[label=\alph*.,leftmargin=2em]
  \item Fill in the table above showing the payments debt and equity holders of each firm will receive
        given each of the two possible levels of free cash flows.
  \item Suppose you hold 10\% of the equity of ABC. What is another portfolio you could hold that would
        provide the same cash flows?
  \item Suppose you hold 10\% of the equity of XYZ. If you can borrow at 10\%, what is an alternative
        strategy that would provide the same cash flows?
\end{enumerate}}

% ---------------------------------------------------------------- Question 2.3
\newcommand{\QtwoC}{%
\srcline{Berk, DeMarzo \& Harford, \emph{Fundamentals of Corporate Finance} 2e, Ch.~13, Exercise 13.21,
p.~405 (PDF 436) $\cdot$ Answer in book: no}

\noindent Your company has two divisions: One division sells software and the other division sells
computers through a direct sales channel, primarily taking orders over the Internet. You have decided
that Dell Computer is very similar to your computer division, in terms of both risk and financing. You
go online and find the following information: Dell's beta is 1.21, the risk-free rate is 4.5\%, its
market value of equity is \$67 billion, and it has \$700 million worth of debt with a yield to maturity
of 6\%. Your tax rate is 35\% and you use a market risk premium of 5\% in your WACC estimates.
\begin{enumerate}[label=\alph*.,leftmargin=2em]
  \item What is an estimate of the WACC for your computer sales division?
  \item If your overall company WACC is 12\% and the computer sales division represents 40\% of the
        value of your firm, what is an estimate of the WACC for your software division?
\end{enumerate}}

% ---------------------------------------------------------------- Question 2.4
\newcommand{\QtwoD}{%
\srcline{Berk, DeMarzo \& Harford, \emph{Fundamentals of Corporate Finance} 2e, Ch.~16, Exercise 16.6,
p.~492 (PDF 523) $\cdot$ Answer in book: no}

\noindent Suppose Microsoft has no debt and a WACC of 9.2\%. The average debt-to-value ratio for the
software industry is 5\%. What would its cost of equity be if it took on the average amount of debt for
its industry at a cost of debt of 6\%?}

% ---------------------------------------------------------------- Question 3.1
\newcommand{\QthreeA}{%
\srcline{Berk, DeMarzo \& Harford, \emph{Fundamentals of Corporate Finance} 2e, Ch.~16, Exercises
16.7--16.8, p.~492 (PDF 523) $\cdot$ Answer in book: no}

\noindent\textbf{7.} Pelamed Pharmaceuticals had EBIT of \$325 million in 2010. In addition, Pelamed had
interest expenses of \$125 million and a corporate tax rate of 40\%.
\begin{enumerate}[label=\alph*.,leftmargin=2em]
  \item What was Pelamed's 2010 net income?
  \item What was the total of Pelamed's 2010 net income and interest payments?
  \item If Pelamed had no interest expenses, what would its 2010 net income have been? How does it
        compare to your answer in part (b)?
  \item What was the amount of Pelamed's interest tax shield in 2010?
\end{enumerate}
\noindent\textbf{8.} Grommit Engineering expects to have net income next year of \$20.75 million and free
cash flow of \$22.15 million. Grommit's marginal corporate tax rate is 35\%.
\begin{enumerate}[label=\alph*.,leftmargin=2em]
  \item If Grommit increases leverage so that its interest expense rises by \$1 million, how will its
        net income change?
  \item For the same increase in interest expense, how will its free cash flow change?
\end{enumerate}}

% ---------------------------------------------------------------- Question 3.2
\newcommand{\QthreeB}{%
\srcline{Berk, DeMarzo \& Harford, \emph{Fundamentals of Corporate Finance} 2e, Ch.~16, Exercises 16.9
and 16.12, p.~492 (PDF 523) $\cdot$ Answer in book: no}

\noindent\textbf{9.} Assume that Microsoft has a total market value of \$300 billion and a marginal tax
rate of 35\%. If it permanently changes its leverage from no debt by taking on new debt in the amount of
13\% of its current market value, what is the present value of the tax shield it will create?

\smallskip\noindent\textbf{12.} Arnell Industries has \$10 million in permanent debt outstanding. The
firm will pay interest only on this debt. Arnell's marginal tax rate is expected to be 35\% for the
foreseeable future.
\begin{enumerate}[label=\alph*.,leftmargin=2em]
  \item Suppose Arnell pays interest of 6\% per year on its debt. What is its annual interest tax
        shield?
  \item What is the present value of the interest tax shield, assuming its risk is the same as the loan?
  \item Suppose instead that the interest rate on the debt is 5\%. What is the present value of the
        interest tax shield in this case?
\end{enumerate}}

% ---------------------------------------------------------------- Question 3.3
\newcommand{\QthreeC}{%
\srcline{Berk, DeMarzo \& Harford, \emph{Fundamentals of Corporate Finance} 2e, Ch.~16, Exercises
16.13--16.15, p.~493 (PDF 524) $\cdot$ Answer in book: no}

\noindent\textbf{13.} Rogot Instruments makes fine violins, violas, and cellos. It has \$1 million in debt
outstanding, equity valued at \$2 million, and pays corporate income tax at a rate of 35\%. Its cost of
equity is 12\% and its cost of debt is 7\%.
\begin{enumerate}[label=\alph*.,leftmargin=2em]
  \item What is Rogot's pretax WACC?
  \item What is Rogot's (effective after-tax) WACC?
\end{enumerate}
\noindent\textbf{14.} Rumolt Motors has 30 million shares outstanding with a price of \$15 per share. In
addition, Rumolt has issued bonds with a total current market value of \$150 million. Suppose Rumolt's
equity cost of capital is 10\%, and its debt cost of capital is 5\%.
\begin{enumerate}[label=\alph*.,leftmargin=2em]
  \item What is Rumolt's pretax WACC?
  \item If Rumolt's corporate tax rate is 35\%, what is its after-tax WACC?
\end{enumerate}
\noindent\textbf{15.} Summit Builders has a market debt-equity ratio of 0.65 and a corporate tax rate of
40\%, and it pays 7\% interest on its debt. By what amount does the interest tax shield from its debt
lower Summit's WACC?}

% ---------------------------------------------------------------- Question 3.4
\newcommand{\QthreeD}{%
\srcline{Berk, DeMarzo \& Harford, \emph{Fundamentals of Corporate Finance} 2e, Ch.~17, Exercises
17.3--17.5, p.~527 (PDF 558) $\cdot$ Answer in book: no. Problem 5 refers to Problem 4, so both are in
this block.}

\noindent\textbf{3.} ECB Co. has 1 million shares outstanding selling at \$20 per share. It plans to
repurchase 100,000 shares at the market price. What will its market capitalization be after the
repurchase? What will its stock price be?

\smallskip\noindent\textbf{4.} KMS corporation has assets of \$500 million, \$50 million of which are
cash. It has debt of \$200 million. If KMS repurchases \$20 million of its stock:
\begin{enumerate}[label=\alph*.,leftmargin=2em]
  \item What changes will occur on its balance sheet?
  \item What will its new leverage ratio be?
\end{enumerate}
\noindent\textbf{5.} Suppose that KMS in Problem 4 decides to initiate a dividend instead, but it wants
the present value of the payout to be the same \$20 million. If its cost of equity capital is 10\%, to
what amount per year in perpetuity should it commit (assuming perfect capital markets)?}


\begin{document}

\begin{center}
  {\Large\bfseries Finance II --- Mock Midterm Exam, Set 3 (Topics 1, 2A, 2B)}\\[4pt]
  {\large Worked solutions}
\end{center}

{\sloppy\noindent Exam paper: \texttt{Simulados/mock3-midterm-topics-1-2b.pdf}. Each question's full
statement (grey box, with its source) comes first, then the solution, a boxed answer and a marking guide.
True/false items earn 1 mark for the verdict and 1 for the reason; each verdict is compared with the
book's printed answer where one exists. Higgins Ch.~2, Exercise 3 (Part~1.1) also appears in Lista R1.
Every number is reproduced by \texttt{Simulados/mock3-midterm-topics-1-2b-solutions\_check.py}.
References: lecture slides \texttt{Aula\_Topic1}, \texttt{Aula\_Topic2A}, \texttt{Aula\_Topic2B} and
the formula sheet \texttt{Notas/formula-sheet-topics-1-2b.pdf}.\par}

% =============================================================================
\section{Question 1.1 --- Financial ratios: true or false}

\statement{\QoneA}

\begin{itemize}[leftmargin=2em]
  \tf{a.}{True}{$A=L+E$, so $A/E=(L+E)/E=1+L/E$. This is an identity, and it is the leverage factor in
  the DuPont decomposition $ROE=\frac{NI}{S}\cdot\frac SA\cdot\frac AE$. \emph{Book: true.}}
  \tf{b.}{True}{$ROE=NI/E$ and $ROA=NI/A$ share the numerator, and $A\ge E$ whenever liabilities are
  non-negative, so $ROE=ROA\cdot\frac AE\ge ROA$. Caveat worth stating in an exam: with a \emph{loss}
  ($NI<0$) leverage magnifies it and $ROE<ROA$; ``always'' holds only for non-negative income.
  \emph{Book: true.}}
  \tf{c.}{False}{A payables period longer than the collection period is convenient (suppliers then
  finance the receivables), but both are set by industry practice and bargaining power; nothing says
  it \emph{should always} hold. \emph{Book: false.}}
  \tf{d.}{True}{The acid-test ratio is $(CA-\text{inventory})/CL$; inventory cannot be negative, so the
  numerator is never larger than the current ratio's. \emph{Book: true.}}
  \tf{e.}{True}{Higher turnover means the same sales on fewer assets: with margin and leverage fixed,
  ROA and ROE rise (DuPont). ``All else equal'' is what makes it true. \emph{Book: true.}}
  \tf{f.}{False}{The earnings yield is $E/P$, the inverse of $P/E$; equal earnings yields imply equal
  P/E ratios. \emph{Book: false.}}
  \tf{g.}{False}{Without taxes and transaction costs an unrealized gain can be turned into cash by
  selling, so it is worth exactly as much as a realized gain. Value is what the asset can be sold for,
  not what has been booked. \emph{Book: false.}}
\end{itemize}

\answer{a T, b T (for non-negative income), c F, d T, e T, f F, g F.}
\refline{Topic 1 slides (balance sheet, DuPont, ratios); Higgins Ch.~2.}
\rubric{14}{per item: 1 for the verdict, 1 for a correct reason (the identity, the shared numerator, the
inverse relation, etc.).}

% =============================================================================
\section{Question 1.2 --- Growth, new equity and repurchases: true or false}

\statement{\QoneB}

\begin{itemize}[leftmargin=2em]
  \tf{a.}{False}{The sustainable growth rate is $g^*=P\cdot R\cdot A\cdot\hat T$ (profit margin,
  retention, asset turnover, assets-to-equity). A firm can exceed it by raising any of the four ratios,
  not only by issuing stock; the limit is how far each ratio can be pushed. \emph{Book: false.}}
  \tf{b.}{False}{Firms with heavy losses usually find equity expensive or unavailable: outside investors
  doubt the cash flows. Exceptions are ``glamorous'' firms with an exciting story; traditional loss-makers
  struggle. \emph{Book: false.}}
  \tf{c.}{True}{Fewer shares raise EPS, while the cash spent (lost interest, or new interest if
  borrowed) lowers earnings; usually the first effect wins. But in a perfect market the higher EPS is
  riskier: the price does not rise and the P/E falls (MM). \emph{Book: true.}}
  \tf{d.}{True}{Survey evidence: managers usually think their shares are undervalued, and buying
  undervalued shares transfers value from selling to remaining shareholders. (This is an information
  effect, outside the perfect-market benchmark.) \emph{Book: true.}}
  \tf{e.}{False}{Slow-growth firms have growth problems too: they generate more cash than they can
  invest well and face pressure to grow through value-destroying investments. \emph{Book: false.}}
  \tf{f.}{False}{Only growth earning more than the cost of capital (positive NPV) raises the price;
  growth at returns below the cost of capital destroys value. \emph{Book: false.}}
\end{itemize}

\answer{a F, b F, c T, d T, e F, f F.}
\refline{Topic 1 slides (new equity financing); Topic 2A (NPV rule; EPS vs.\ value under MM); Higgins
Ch.~4.}
\rubric{12}{per item: 1 for the verdict, 1 for the reason. In (c) full marks need the EPS mechanism; the
MM remark is a plus.}

% =============================================================================
\section{Question 1.3 --- Risk-free debt and the cost of equity}

\statement{\QoneC}

The argument is \textbf{wrong}. Equity is the residual claim. Risk-free debt takes a fixed payment, so all
the variation of the asset cash flows now falls on a smaller equity stake:
\[
  \beta_E=\beta_A+\frac DE(\beta_A-\beta_D)=\beta_A\Big(1+\frac DE\Big)>\beta_A\quad(\beta_D=0).
\]
Equity becomes riskier and its cost rises, $r_E=r_A+\frac DE(r_A-r_D)$, by exactly enough to keep the
WACC equal to $r_A$ (MM II). The ``cheap'' debt is offset by more expensive equity; without taxes, the
cost of capital of the firm does not change.

\answer{False: risk-free debt raises the equity's risk and cost ($\beta_E=\beta_A(1+D/E)$); the WACC stays
at $r_A$.}
\rubric{2}{1 for ``equity risk rises'', 1 for ``WACC unchanged (MM II)''.}

% =============================================================================
\section{Question 2.1 --- Biotech start-up: MM and the fraction of equity sold}

\statement{\QtwoIntro\par\medskip\QtwoA}

\step{1}{Value without leverage}
\$2 million buys 50\% of the unlevered equity, so the whole firm is worth $V_U=2/0.5=\$4$ million. (A
\$30 million payoff in the success state is worth 4 today, so the price of \$1 paid on success is
$4/30=0.1333$.)

\step{2}{Borrow \$1 million}
MM I: $V_L=V_U=4$. With debt worth 1, the equity is worth $E=4-1=3$. To raise the remaining \$1
million, sell $1/3=33.33\%$ of the equity. The debt is risky (it gets nothing if the research fails),
but MM holds for any debt: the claims just split the same \$4 million.

\step{3}{Your stake}
(a) $50\%\times4=\$2$ million; (b) $\frac23\times3=\$2$ million.

\answer{(a) \$4 million. (b) Sell one third of the equity. (c) \$2 million in both cases: the financing
choice does not change your wealth.}
\intuition{Debt only changes how the \$4 million is split; outside investors pay a fair price for what
they get, so the founder's wealth is the same.}
\refline{Topic 2A slides (MM I, perfect markets); Berk Ch.~16.1.}
\rubric{8}{2 for $V_U=4$; 3 for $E=3$ and the fraction 1/3; 3 for both stakes worth 2 and the reason.}

% =============================================================================
\section{Question 2.2 --- Homemade leverage}

\statement{\QtwoB}

\step{1}{Payments (part a)}
XYZ pays interest $0.10\times5000=\$500$ in either state.
\begin{center}
\begin{tabular}{lcccc}
  \toprule
  & \multicolumn{2}{c}{ABC} & \multicolumn{2}{c}{XYZ} \\
  \cmidrule(lr){2-3}\cmidrule(lr){4-5}
  FCF & Debt Payments & Equity Dividends & Debt Payments & Equity Dividends \\ \midrule
  \$\ 800 & 0 & 800 & 500 & 300 \\
  \$1000 & 0 & 1000 & 500 & 500 \\ \bottomrule
\end{tabular}
\end{center}

\step{2}{Unlevering (part b)}
10\% of ABC's equity pays 80 or 100. Hold 10\% of XYZ's equity \emph{and} 10\% of XYZ's debt
($0.1\times5000=\$500$ of debt): $30+50=80$ or $50+50=100$. Same cash flows.

\step{3}{Homemade leverage (part c)}
10\% of XYZ's equity pays 30 or 50. Borrow \$500 at 10\% (interest \$50) and buy 10\% of ABC's equity:
$80-50=30$ or $100-50=50$. Same cash flows.

\answer{(a) Table above. (b) 10\% of XYZ's equity plus 10\% (\$500) of XYZ's debt. (c) Borrow \$500 at
10\% and buy 10\% of ABC's equity.}
\intuition{Investors can create or undo leverage on their own account at the same rate as firms, so
they will not pay more for a levered or an unlevered firm: by no arbitrage, $V_L=V_U$ (MM I).}
\refline{Topic 2A slides (homemade leverage, MM I proof); Berk Ch.~16.1.}
\rubric{10}{4 for the table; 3 for (b); 3 for (c), each with the payoffs shown in both states.}

% =============================================================================
\section{Question 2.3 --- Divisional WACC from a comparable}

\statement{\QtwoC}

\step{1}{Comparable's cost of equity}
$r_E=4.5\%+1.21\times5\%=10.55\%$.

\step{2}{WACC of the computer division}
Dell matches the division in risk \emph{and} financing, so its WACC can be used directly (no need to
unlever and relever). With $V=67+0.7=67.7$, $E/V=0.9897$ and $D/V=0.0103$; the after-tax cost of debt is
$6\%(1-0.35)=3.9\%$:
\[
  r_{WACC}=0.9897(10.55\%)+0.0103(3.9\%)=10.48\%.
\]

\step{3}{WACC of the software division}
The firm's WACC is the value-weighted average of the divisions' WACCs:
$12\%=0.40(10.48\%)+0.60\,r_{soft}$, so $r_{soft}=(12\%-0.40\times10.48\%)/0.60=13.01\%$.

\answer{(a) About 10.48\%. (b) About 13.01\%.}
\intuition{Each division is discounted at its own risk; using the firm-wide 12\% would over-penalise the
computer division and favour the riskier software division.}
\refline{Topic 2A slides (project cost of capital, comparables); Berk Ch.~13.}
\rubric{8}{2 for $r_E$; 3 for the computer-division WACC with market weights and after-tax $r_D$; 3 for
the software WACC from the weighted average.}

% =============================================================================
\section{Question 2.4 --- MM II: Microsoft with industry leverage}

\statement{\QtwoIntro\par\medskip\QtwoD}

With no debt, the WACC is the unlevered cost of capital: $r_U=9.2\%$. The section assumes no taxes, so MM
II applies with $D/E=0.05/0.95=0.0526$:
\[
  r_E=r_U+\frac DE(r_U-r_D)=9.2\%+0.0526(9.2\%-6\%)=9.37\%.
\]
Check: $0.95(9.37\%)+0.05(6\%)=9.2\%$, so the WACC is unchanged.

\answer{About 9.37\%.}
\rubric{6}{2 for $r_U=9.2\%$; 2 for $D/E=0.0526$ (not 0.05); 2 for $r_E$ and the WACC check.}

% =============================================================================
\section{Question 3.1 --- Interest, net income and free cash flow}

\statement{\QthreeA}

\step{1}{Pelamed (Problem 7)}
(a) $NI=(325-125)(1-0.40)=\$120$ million. (b) $NI+\text{interest}=120+125=\$245$ million. (c) Without
interest, $NI=325(0.60)=\$195$ million, which is \$50 million \emph{less} than in (b). (d) The interest
tax shield is $0.40\times125=\$50$ million: exactly the gap between (b) and (c).

\step{2}{Grommit (Problem 8)}
(a) An extra \$1 million of interest lowers taxable income by 1 and taxes by 0.35, so net income falls by
$1(1-0.35)=\$0.65$ million, to \$20.10 million. (b) Free cash flow is unlevered (computed before
interest), so it does not change: \$22.15 million.

\answer{7: (a) \$120m; (b) \$245m; (c) \$195m, \$50m less than (b); (d) \$50m. 8: (a) falls by \$0.65m
(to \$20.10m); (b) unchanged.}
\intuition{Debt and equity together receive more when the firm pays interest, because the government
receives less: that extra \$50 million is the tax shield. It never enters FCF; it enters valuation through
the after-tax WACC or APV.}
\refline{Topic 2B slides (interest tax shield); Berk Ch.~16.3.}
\rubric{10}{6 for Problem 7 (1.5 per part, with the link (b)$-$(c)$=$(d)); 4 for Problem 8 (2 each; (b)
needs the reason ``FCF is unlevered'').}

% =============================================================================
\section{Question 3.2 --- Present value of permanent tax shields}

\statement{\QthreeB}

\step{1}{Microsoft (Problem 9)}
Permanent debt of $0.13\times300=\$39$ billion; its shields are as safe as the debt, so
$PV(ITS)=\tau D=0.35\times39=\$13.65$ billion.

\step{2}{Arnell (Problem 12)}
(a) $ITS=0.35\times0.06\times10=\$0.21$ million a year. (b) A perpetuity discounted at the debt's rate:
$0.21/0.06=\$3.5$ million. (c) At 5\%: $ITS=\$0.175$ million and $PV=0.175/0.05=\$3.5$ million.

\answer{9: \$13.65 billion. 12: (a) \$0.21m; (b) \$3.5m; (c) \$3.5m.}
\intuition{For permanent debt, $PV(ITS)=\tau r_DD/r_D=\tau D$: the rate cancels, because a higher rate
raises the shield and the discount rate in the same proportion.}
\refline{Topic 2B slides (debt-policy table, fixed perpetual debt); Berk Ch.~16.3.}
\rubric{8}{2 for Problem 9; 2 for 12(a); 2 for 12(b); 2 for 12(c) with the reason it equals $\tau D$.}

% =============================================================================
\section{Question 3.3 --- Pretax and after-tax WACC}

\statement{\QthreeC}

\step{1}{Rogot (Problem 13)}
$V=1+2=3$, so $E/V=2/3$, $D/V=1/3$. (a) Pretax: $\frac23(12\%)+\frac13(7\%)=10.33\%$. (b) After-tax:
$\frac23(12\%)+\frac13(7\%)(0.65)=9.52\%$.

\step{2}{Rumolt (Problem 14)}
$E=30\times15=\$450$ million, $D=150$, $V=600$. (a) $0.75(10\%)+0.25(5\%)=8.75\%$. (b)
$0.75(10\%)+0.25(5\%)(0.65)=8.31\%$.

\step{3}{Summit (Problem 15)}
$D/V=0.65/1.65=0.3939$. The shield lowers the WACC by $\tau r_D\frac DV=0.40(7\%)(0.3939)=1.10$
percentage points.

\answer{13: (a) 10.33\%; (b) 9.52\%. 14: (a) 8.75\%; (b) 8.31\%. 15: by about 1.10 percentage points.}
\intuition{The pretax WACC is the return on the firm's assets ($r_U$ when $D/V$ is held constant). The
after-tax WACC is lower by $\tau r_D D/V$: the tax shield, expressed as a rate.}
\refline{Topic 2B slides (WACC $=r_U-\tau r_D D/V$); Berk Ch.~16.3.}
\rubric{12}{4 for Rogot; 4 for Rumolt (market value of equity first); 4 for Summit (2 for converting
$D/E$ to $D/V$, 2 for the reduction).}

% =============================================================================
\section{Question 3.4 --- Repurchases and dividends in a perfect market}

\statement{\QthreeD}

\step{1}{ECB (Problem 3)}
Market cap before: $1\times20=\$20$ million. The firm spends $0.1\times20=\$2$ million, so the market
cap falls to \$18 million and the price stays at $18/0.9=\$20$: buying at a fair price is zero-NPV.

\step{2}{KMS (Problem 4)}
(a) Cash falls from 50 to 30 and assets from 500 to 480; debt stays at 200; equity falls from 300 to
280. (b) $D/E$ rises from $200/300=0.667$ to $200/280=0.714$ (equivalently $D/V$ from 40\% to 41.7\%).

\step{3}{KMS dividend (Problem 5)}
A perpetuity worth \$20 million at 10\% pays $C=20\times0.10=\$2$ million a year.

\answer{3: \$18 million, \$20 a share. 4: (a) cash and assets $-20$, equity $-20$ (to 280); (b) $D/E$
from 0.667 to 0.714. 5: \$2 million a year.}
\intuition{Paying out cash shrinks the firm but not the value per share in a perfect market; using cash
(negative debt) to buy back shares raises leverage.}
\refline{Topic 2B slides (payout irrelevance); Berk Ch.~17.2.}
\rubric{10}{3 for Problem 3 (cap and price, with the reason); 4 for Problem 4; 3 for Problem 5.}

\end{document}
